\documentclass[border=0pt]{standalone}
% The shared preamble of every docs diagram (docs/diagrams/<name>.tex).
% dvisvgm's driver first, so patterns and opacity survive into the SVG.
\def\pgfsysdriver{pgfsys-dvisvgm.def}
\usepackage[siunitx,RPvoltages]{circuitikz}
\usepackage{xcolor}
% Latin Modern: outline fonts for every glyph a label uses, µ included.
\usepackage[T1]{fontenc}
\usepackage{lmodern}
\usetikzlibrary{arrows.meta,backgrounds,calc,patterns,positioning}
\tikzset{
  % A white card, so diagrams read in mdBook's dark themes too.
  every picture/.append style={show background rectangle,
    background rectangle/.style={fill=white,inner sep=4pt}},
  % Layout geometry: layers told apart by fill and hatching, never colour alone.
  layer a/.style={draw=black,line width=0.6pt,fill=gray!20},
  layer b/.style={draw=black,line width=0.6pt,fill=gray!45},
  layer c/.style={draw=black,line width=0.6pt,fill=white,postaction={pattern=north east lines}},
  layer d/.style={draw=black,line width=0.6pt,fill=white,postaction={pattern=crosshatch dots}},
  marker/.style={draw=black,dashed,line width=0.6pt},
  % A measured distance and its label.
  dim/.style={{Latex[length=2mm]}-{Latex[length=2mm]},line width=0.5pt},
  dim label/.style={font=\small,fill=white,inner sep=1pt},
  % Where a rule fails: thick, dashed and labelled, not only red.
  violation/.style={draw=red!75!black,line width=1.4pt,dash pattern=on 3pt off 2pt},
  violation label/.style={font=\small\bfseries,text=red!75!black},
  caption/.style={font=\small},
}

\begin{document}
\begin{circuitikz}
  % A CMOS inverter at x = #1 between rails at y = 2.6 and -2.6; input at
  % (#1-0.9, 0), output at (#1, 0).
  \newcommand{\inv}[2]{%
    \draw (#1,1) node[pmos,anchor=D](#2p){} (#1,-1) node[nmos,anchor=D](#2n){};
    \draw (#2p.D) -- (#2n.D);
    \draw (#2p.G) -- (#2n.G) coordinate[pos=0.5](#2in);
    \draw (#2p.S) -- (#2p.S |- 0,2.6);
    \draw (#2n.S) -- (#2n.S |- 0,-2.6);
  }
  \inv{0}{a}
  \inv{4.5}{b}
  \draw (-1.6,2.6) -- (6,2.6) node[right]{VDD, 3.3 V};
  \draw (-1.6,-2.6) -- (6,-2.6) node[right]{VSS, 0 V};
  \draw (ain) to[short,*-] (-1.6,0) node[left]{A};
  \draw (0,0) to[short,*-*] (bin);
  \draw (4.5,0) to[short,*-] (5.5,0) node[right]{Y: 0 V to 3.3 V};
  \node[caption,align=center,anchor=north] at (-2.2,-0.3) {no supply:\\not examined};
  \node[caption,anchor=south] at (1.8,0.05) {Y1: 0 V to 3.3 V};
  \draw[violation] (3.45,-2.1) rectangle (5.0,2.1);
  \node[violation label,anchor=north] at (4.25,-2.75) {gate oxide 3.3 V $>$ 1.95 V};
\end{circuitikz}
\end{document}
